\section{Estructura de la célula procariota}

Las células procariotas son organismos unicelulares que carecen de núcleo definido y organelos membranosos. Su material genético se encuentra disperso en el citoplasma en una región llamada nucleoide.

\section{Características principales}

\begin{itemize}
    \item Tamaño: 1-10 µm
    \item Sin núcleo definido (nucleoide)
    \item Sin organelos membranosos (mitocondrias, cloroplastos, retículo endoplásmico)
    \item Ribosomas 70S (menores que los eucariotas)
    \item Pared celular compuesta por peptidoglicano
    \item Reproducción asexual por fisión binaria
\end{itemize}

\section{Estructura de la célula bacteriana}

\subsection{Pared celular}
\begin{itemize}
    \item Bacterias Gram positivas: Capa gruesa de peptidoglicano
    \item Bacterias Gram negativas: Capa fina de peptidoglicano + membrana externa
\end{itemize}

\subsection{Membrana plasmática}
\begin{itemize}
    \item Bicapa lipídica similar a la de eucariotas
    \item Función: Barrera selectiva, transporte, respiración
\end{itemize}

\subsection{Citoplasma}
\begin{itemize}
    \item Contiene ribosomas 70S
    \item Gránulos de almacenamiento
    \item Plasmidios (ADN extracromosómico)
\end{itemize}

\subsection{Nucleoide}
\begin{itemize}
    \item Región que contiene el ADN circular
    \item Sin membrana nuclear
\end{itemize}

\subsection{Flagelos y fimbrias}
\begin{itemize}
    \item Flagelos: Para movilidad
    \item Fimbrias: Para adhesión
    \item Pili: Para conjugación
\end{itemize}

\section{Diversidad de procariotas}

\begin{itemize}
    \item Bacterias: Eubacterias (E. coli, Bacillus)
    \item Arqueas: Extremófilas (ambientes extremos)
    \item Cianobacterias: Fotosintéticas
\end{itemize}

\subsection{Aplicaciones biotecnológicas}
\begin{itemize}
    \item Producción de antibióticos
    \item Fermentación industrial
    \item Producción de enzimas
    \item Biorremediación
\end{itemize}